\documentclass[11pt,a4paper,parskip=half-]{scrartcl}
\usepackage[T1]{fontenc}
\usepackage[utf8]{inputenc}
\usepackage{mathpazo}               % Palatino text and maths
\usepackage[scaled=0.9]{helvet}     % sans-serif for headings, scaled to match
\usepackage[a4paper,margin=2.5cm]{geometry}
\usepackage{microtype}
\usepackage{amsmath}
\usepackage{booktabs}
\usepackage{tabularx}
\usepackage{enumitem}
\setlist{nosep}
\usepackage{xcolor}
\usepackage{siunitx}
\sisetup{group-minimum-digits=4}   % 4000 and 52000 are both written with a thin space
\usepackage{pgfgantt}               % the Gantt chart; it loads TikZ
\usepackage[round,authoryear]{natbib}
\bibliographystyle{plainnat}
\usepackage[colorlinks=true,linkcolor=black,citecolor=blue!50!black,urlcolor=blue!50!black]{hyperref}
\usepackage[capitalise,noabbrev]{cleveref}
\addtokomafont{disposition}{\sffamily}

% ---- Fill in once -------------------------------------------------------------
\newcommand{\proposaltitle}{Title of the Research Proposal}
\newcommand{\applicant}{First Applicant}
\newcommand{\institution}{Example University, Department of Science}
\newcommand{\projectduration}{24 months}
\hypersetup{pdftitle={\proposaltitle}, pdfauthor={\applicant}}

\usepackage[headsepline]{scrlayer-scrpage}
\clearpairofpagestyles
\ihead{\small\proposaltitle}
\ohead{\small\applicant}
\cfoot{\pagemark}

\begin{document}

% ---- Cover block: the facts a reviewer looks for first ---------------------------
{\raggedright
\sffamily\bfseries\LARGE \proposaltitle\par}
\medskip
\noindent\textsf{\large Research proposal}\par
\vspace{0.6em}
\noindent\begin{tabularx}{\textwidth}{@{}l X@{}}
  \toprule
  \sffamily Applicant       & \applicant, \institution \\
  \sffamily Duration        & \projectduration \\
  \sffamily Funding sought  & see \cref{sec:budget} \\
  \sffamily Keywords        & keyword one, keyword two, keyword three \\
  \bottomrule
\end{tabularx}

\section*{Summary}
One paragraph, readable by a non-specialist, of about 150 words: the problem, the approach,
what will exist at the end that does not exist now, and why it matters. Panels often read only
this and the aims; write them first and with the most care.

% ---- 1 -----------------------------------------------------------------------------
\section{Aims and objectives}\label{sec:aims}
\textbf{Overall aim.} One sentence: what this project will achieve.

\textbf{Objectives.} Three to five, each one checkable at the end of the project. Number them
so that the work plan, the milestones and the budget can refer to them.
\begin{description}[leftmargin=2.2em,labelwidth=2em,labelsep=0.2em,font=\bfseries]
  \item[O1] Establish how X behaves under condition Y.
  \item[O2] Develop a method that does Z with at least half the error of the current best approach.
  \item[O3] Validate the method on real data and release it as open software.
\end{description}

% ---- 2 -----------------------------------------------------------------------------
\section{Background and state of the art}
Show that you know the field and that the gap is real. Cite the work a reviewer expects to
see, say what it achieves and where it stops, and end with the question this project answers.
Data should follow the FAIR principles \citep{wilkinson2016}; planned work is scheduled
as in a Gantt chart \citep{gantt1919}. \citet{kingma2015} is an example of a citation in running text.

% ---- 3 -----------------------------------------------------------------------------
\section{Methods}
Say how each objective will be reached, in the same order as in \cref{sec:aims}. Give
enough detail that a specialist can judge feasibility: data, equipment, analysis, and how
you will know it worked.
\begin{equation}
  \hat{\theta} = \arg\min_{\theta}\; \frac{1}{n}\sum_{i=1}^{n} \ell\bigl(f_{\theta}(x_i),\, y_i\bigr)
  \label{eq:objective}
\end{equation}
Equations are numbered as usual (\cref{eq:objective}). Name the main risks and what you
would do if they happen, either here or in a short table.

% ---- 4 -----------------------------------------------------------------------------
\section{Work plan}\label{sec:plan}
The work is divided into work packages (WP), each with a start and end month, an
objective it serves, and a lead. \Cref{fig:gantt} shows the schedule.

\begin{description}[leftmargin=2.8em,labelwidth=2.6em,labelsep=0.2em,font=\bfseries]
  \item[WP1] \textbf{Preparation} (months 1--6, serves O1). Literature review, ethics approval, data collection set-up.
  \item[WP2] \textbf{Method development} (months 4--15, serves O2). Design, implementation and tests on synthetic data.
  \item[WP3] \textbf{Validation} (months 13--21, serves O3). Experiments on real data and comparison with existing methods.
  \item[WP4] \textbf{Dissemination} (months 18--24, serves O3). Papers, software release, final report.
\end{description}

% ================================================================================
% HOW TO EDIT THE GANTT CHART
%   * The chart runs from month 1 to month 24: change the two numbers in
%     \begin{ganttchart}{1}{24}, the \gantttitle lines (one per year) and the list
%     in \gantttitlelist if your project is longer or shorter.
%   * One line per task. Copy a line and change the label and the two month numbers:
%         \ganttbar{label}{first month}{last month} \\
%     \ganttgroup draws a bracket-style bar for a whole work package; \ganttmilestone
%     draws a diamond at one month. Every line but the last ends with \\ .
%   * If the chart gets too wide or too narrow, change x unit (width of one month).
% ================================================================================
\begin{figure}[ht]
  \centering
  \begin{ganttchart}[
      x unit=0.44cm, y unit title=0.5cm, y unit chart=0.62cm,
      title label font=\sffamily\tiny, bar label font=\sffamily\footnotesize,
      group label font=\sffamily\footnotesize\bfseries, milestone label font=\sffamily\footnotesize,
      hgrid, vgrid={*{5}{draw=none}, dotted},       % a dotted line every sixth month
      bar/.append style={fill=blue!45!black!70, draw=none},
      group/.append style={fill=black!70, draw=none},
      milestone/.append style={fill=orange!85!black, draw=none},
    ]{1}{24}
    \gantttitle{Year 1}{12} \gantttitle{Year 2}{12} \\
    \gantttitlelist{1,...,24}{1} \\
    \ganttgroup{WP1 Preparation}{1}{6} \\
    \ganttbar{Literature review}{1}{4} \\
    \ganttbar{Ethics approval}{2}{6} \\
    \ganttgroup{WP2 Method development}{4}{15} \\
    \ganttbar{Design}{4}{9} \\
    \ganttbar{Implementation}{8}{15} \\
    \ganttgroup{WP3 Validation}{13}{21} \\
    \ganttbar{Experiments}{13}{19} \\
    \ganttbar{Comparison}{18}{21} \\
    \ganttgroup{WP4 Dissemination}{18}{24} \\
    \ganttbar{Papers and software}{18}{24} \\
    \ganttmilestone{Milestones}{6} \ganttmilestone{}{15} \ganttmilestone{}{21} \ganttmilestone{}{24}
  \end{ganttchart}
  \caption{Work plan. Diamonds are the milestones M1--M4 of \cref{tab:milestones}.}
  \label{fig:gantt}
\end{figure}

% ---- 5 -----------------------------------------------------------------------------
\section{Milestones and deliverables}
A milestone is a point at which something can be checked. Keep the months in step with
the diamonds in \cref{fig:gantt}.

\begin{table}[ht]
  \centering
  \caption{Milestones (M) and deliverables (D).}
  \label{tab:milestones}
  \begin{tabularx}{\textwidth}{@{}l c l X@{}}
    \toprule
    ID & Month & WP & Description and how it is verified \\
    \midrule
    M1 & 6  & WP1 & Data collection set up and approved; checked by the ethics committee. \\
    M2 & 15 & WP2 & Method implemented and passing all tests on synthetic data. \\
    D1 & 15 & WP2 & Open-source software, version 1.0, with documentation. \\
    M3 & 21 & WP3 & Validation complete; results meet the target in O2. \\
    D2 & 24 & WP4 & Two journal papers submitted; final report to the funder. \\
    M4 & 24 & WP4 & Project closed and data archived under FAIR principles. \\
    \bottomrule
  \end{tabularx}
\end{table}

% ---- 6 -----------------------------------------------------------------------------
\section{Budget}\label{sec:budget}
% ================================================================================
% HOW TO EDIT THE BUDGET
%   One line per cost item:   \budgetline{description}{year 1 amount}{year 2 amount}
%   Whole amounts in euros (no decimals, no thousands separator). Copy, delete or reorder
%   lines freely. The sums, the overhead and the total are calculated for you.
%   Change the overhead rate in \budgetoverhead{15} (percent of direct costs) and the
%   currency in the caption. For three years, add a counter and a column in the same way.
% ================================================================================
\newcounter{bA}\newcounter{bB}   % direct costs, year 1 and year 2
\newcounter{oA}\newcounter{oB}   % overhead, year 1 and year 2
% siunitx does not calculate, so the sum is evaluated first and handed over as plain digits.
\newcommand{\amount}[1]{\edef\tmpamount{\number\numexpr#1\relax}\expandafter\num\expandafter{\tmpamount}}
\newcommand{\budgetline}[3]{%
  \addtocounter{bA}{#2}\addtocounter{bB}{#3}% add to the running sums
  #1 & \amount{#2} & \amount{#3} & \amount{#2+#3} \\}
\newcommand{\budgetoverhead}[1]{%
  \setcounter{oA}{\numexpr(\value{bA}*#1)/100\relax}%
  \setcounter{oB}{\numexpr(\value{bB}*#1)/100\relax}%
  Indirect costs (#1\,\% of direct) & \amount{\value{oA}} & \amount{\value{oB}}
    & \amount{\value{oA}+\value{oB}} \\}

\begin{table}[ht]
  \centering
  \caption{Requested budget in euros.}
  \label{tab:budget}
  \begin{tabularx}{\textwidth}{@{}X r r r@{}}
    \toprule
    Item & Year 1 & Year 2 & Total \\
    \midrule
    \budgetline{Personnel: one researcher, full time}{52000}{54000}
    \budgetline{Equipment: measurement kit}{8000}{0}
    \budgetline{Consumables}{2000}{2000}
    \budgetline{Travel and conferences}{3000}{4500}
    \budgetline{Publication costs (open access)}{0}{1500}
    \midrule
    Direct costs & \amount{\value{bA}} & \amount{\value{bB}}
      & \amount{\value{bA}+\value{bB}} \\
    \budgetoverhead{15}
    \midrule
    \textbf{Total requested}
      & \textbf{\amount{\value{bA}+\value{oA}}}
      & \textbf{\amount{\value{bB}+\value{oB}}}
      & \textbf{\amount{\value{bA}+\value{bB}+\value{oA}+\value{oB}}} \\
    \bottomrule
  \end{tabularx}
\end{table}

Justify every line above a few thousand euros in a sentence: why it is needed and why it
costs that much. Personnel is normally the largest item and the one reviewers check first.

\bibliography{refs}

\end{document}
